\chapter*{Abstract}
\addcontentsline{toc}{chapter}{Abstract}
We study how spelling errors affect four-topic classification on AG News. The errors swap adjacent interior letters inside eligible words and preserve word length and boundaries. With 20,000 labeled training articles, we compare a word and bigram TF-IDF baseline (\textit{Word TF-IDF LR}) with \textit{Word CNN} and \textit{Subword CNN}. We evaluate clean-text accuracy and accuracy at swap targets up to 20\%. We also test whether replacing half of each topic's training articles with copies edited at the 10\% target improves performance on noisy text.

\textit{Word TF-IDF LR} is the most accurate of the three main models at every swap target. With clean training, its accuracy falls from 90.05\% on clean text to 89.14\% at 20\% swaps. Its median CPU fit time is 13 seconds. The two CNNs fall by about two percentage points, with most of the primary fits' net losses in \emph{Business}.

Training with typos improves \textit{Subword CNN} accuracy by 0.31 to 0.46 percentage points on average across the three 20\% draws. Individual runs vary, and the separate confirmation sample shows a small drop on clean text. A 16-epoch word-only control at seed 42 closes most of the gap between the augmented CNNs at 20\% swaps.

\subsection*{Keywords}
News classification, spelling errors, text representations, convolutional neural networks, data augmentation, robustness evaluation.
